\begin{userprompt}

I need to update my work status on the SEVP portal and I need help to write what my duties are going to be based on the duties that are given on my offer letter

The employee will work as part of an interdisciplinary team of filmmakers, artists, and scientists to produce content focused on scientific topics, applying their background in data science. The employee will be responsible for creating content for the STAGE center’s productions including but not limited to theatre and film. For this purpose, the employee will be testing and comparing AIbased translation tools for content translation and transcription. This work will include evaluating speech-to-text software, segmenting long transcripts into manageable sections, developing structured instructions to improve AI translation quality, and reviewing outputs for accuracy and consistency. In addition, the employee will assist with video editing and the preparation of materials for the productions. Through this role, the employee will gain hands-on experience, receive production credit, develop time and project management skills, and have opportunities for further training in advanced data science techniques relevant to artistic production and science communication.
\end{userprompt}

\begin{botresponse}
Based on your offer letter, here's a concise work status description for the SEVP portal:

\textbf{Work Duties:}

I work as part of an interdisciplinary team producing scientific content for theatre and film productions. My responsibilities include testing and comparing AI-based translation tools for content translation and transcription, evaluating speech-to-text software, segmenting transcripts, developing structured instructions to improve AI translation quality, and reviewing outputs for accuracy and consistency. I also assist with video editing and preparation of production materials, applying my data science background to support the STAGE center's scientific communication projects.
\end{botresponse}

\begin{userprompt}
Is this in line with my major in Applied Data Science ?
\end{userprompt}

\begin{botresponse}
Yes, this is well-aligned with Applied Data Science. Here's why:

\textbf{Core data science elements in your duties:}

\begin{itemize}
\item \textbf{Natural Language Processing (NLP)} : Testing AI translation tools and speech-to-text software
\item \textbf{Model Evaluation} : Comparing different AI tools and reviewing outputs for accuracy and consistency
\item \textbf{Data Processing} : Segmenting long transcripts into manageable sections
\item \textbf{Prompt Engineering} : Developing structured instructions to improve AI translation quality
\item \textbf{Quality Assurance} : Reviewing AI outputs for accuracy and consistency
\end{itemize}

\textbf{Applied aspects:}

\begin{itemize}
\item Real-world application of AI/ML tools in a production environment
\item Working with multimedia data (text, audio, video)
\item Practical problem-solving in scientific communication
\end{itemize}

The role directly applies data science skills to content creation and translation workflows, which is a legitimate application of your field. For SEVP purposes, the key is that you're using AI/ML technologies, data processing, and evaluation methodologies—all core components of applied data science—even though the end product is creative/scientific content rather than traditional tech applications.
\end{botresponse}

\begin{userprompt}
I need another help the legal name of my employer is

Scientists, Technologists, and Artists Generating Exploration, Pritzker School of Molecular Engineering, University of Chicago.

but on the portal there is a character length limit
\end{userprompt}

\begin{botresponse}
Here are some options to fit within character limits:

\textbf{Shortest (most likely to fit):}

\begin{itemize}
\item University of Chicago
\item The University of Chicago
\end{itemize}

\textbf{Medium length:}

\begin{itemize}
\item Pritzker School of Molecular Engineering, University of Chicago
\item STAGE, Pritzker School, University of Chicago
\end{itemize}

\textbf{If more space available:}

\begin{itemize}
\item STAGE Center, Pritzker School of Molecular Engineering, UChicago
\end{itemize}

\textbf{Recommendation:} Start with "University of Chicago" since that's the actual employing institution. The STAGE center and Pritzker School are divisions within the university. If the portal asks for department separately, you can put "STAGE, Pritzker School of Molecular Engineering" there.

Check if your offer letter or employment documents show a specific employer name format—use that if available.
\end{botresponse}

\begin{userprompt}
My offer letter says - The STAGE Center has offer me  the position of .....
\end{userprompt}

\begin{botresponse}
Based on your offer letter stating "The STAGE Center," try these options in order of preference:

\textbf{Best options:}

\begin{enumerate}
\item \textbf{STAGE Center, University of Chicago}
\item \textbf{The STAGE Center, University of Chicago}
\item \textbf{STAGE, Pritzker School, University of Chicago}
\end{enumerate}

\textbf{If extremely limited:}

\begin{itemize}
\item \textbf{University of Chicago} (the legal employing entity)
\end{itemize}

Since your offer letter uses "The STAGE Center" as the hiring unit, option 1 or 2 would be most accurate while keeping it concise. The University of Chicago is still the legal employer, so including that is important for SEVP records.

If there's a separate field for department/division, you could use:

\begin{itemize}
\item \textbf{Employer:} University of Chicago
\item \textbf{Department:} STAGE Center, Pritzker School of Molecular Engineering
\end{itemize}
\end{botresponse}

\begin{userprompt}
since the position is 20 hr / week is it part time ?
\end{userprompt}

\begin{botresponse}
Yes, 20 hours per week is part-time employment.

\textbf{For SEVP portal purposes:}

\begin{itemize}
\item Full-time is typically 40 hours per week
\item Part-time is anything less than full-time
\item 20 hrs/week = part-time position
\end{itemize}

Make sure to accurately report this as part-time on your SEVP portal when updating your work authorization status.
\end{botresponse}

\begin{userprompt}
and it is above the legal required limit right ?
\end{userprompt}

\begin{botresponse}
That depends on your specific visa status and work authorization type:

\textbf{F-1 CPT (Curricular Practical Training):}

\begin{itemize}
\item No minimum hour requirement
\item Can be part-time (20 hrs or less) or full-time (more than 20 hrs)
\item 20 hrs/week is acceptable
\end{itemize}

\textbf{F-1 OPT (Optional Practical Training):}

\begin{itemize}
\item Part-time OPT: minimum 20 hours per week
\item Your 20 hrs/week meets the minimum requirement
\item Must work at least 20 hrs/week to maintain valid OPT status
\end{itemize}

\textbf{On-campus employment (F-1):}

\begin{itemize}
\item Up to 20 hrs/week during academic term
\item Your position meets this limit
\end{itemize}

\textbf{Important:} If you're on OPT, working exactly 20 hours per week is the minimum threshold. If you're on CPT or on-campus employment, 20 hours is fine. Check your I-20 and work authorization documents to confirm which category applies to you, as the requirements differ.
\end{botresponse}

\begin{userprompt}
after graduation
\end{userprompt}

\begin{botresponse}
\textbf{Post-completion OPT (after graduation):}

Yes, 20 hours per week meets the legal minimum requirement.

\textbf{Key rules:}

\begin{itemize}
\item Minimum: 20 hours per week
\item Your position at exactly 20 hrs/week satisfies this requirement
\item You cannot work less than 20 hrs/week or you'll accrue unemployment days
\end{itemize}

\textbf{Important:} On post-completion OPT, you can only have 90 days of unemployment total during your 12-month OPT period. Working part-time (20+ hrs/week) does NOT count as unemployment, so you're compliant.

Make sure your employer reports your actual hours accurately, and keep documentation of your work schedule in case USCIS requests it.
\end{botresponse}

\begin{userprompt}
how do I answer this question Enter the start date for your job with the company that hired you. This date cannot be before your OPT start date. ?
\end{userprompt}

\begin{botresponse}
Enter the actual date you will begin working at the STAGE Center.

\textbf{Important guidelines:}

\begin{enumerate}
\item \textbf{Use your actual first day of work} - the date on your offer letter or when you were told to start
\item \textbf{Must be on or after your OPT start date} - check your I-20 for your OPT start date and ensure your employment start date is not earlier
\item \textbf{Format} - use whatever date format the portal requires (usually MM/DD/YYYY)
\end{enumerate}

\textbf{Example:}

\begin{itemize}
\item If your OPT starts February 1, 2025
\item And your offer letter says you start February 15, 2025
\item Enter: 02/15/2025
\end{itemize}

If your employment start date is somehow before your OPT start date on your I-20, you have a problem and need to contact your DSO (Designated School Official) immediately to resolve it before updating SEVP.
\end{botresponse}

\begin{userprompt}
how is my I 20 going to have my OPT start date ?
\end{userprompt}

\begin{botresponse}
You're right to question this—let me clarify:

\textbf{Your I-20 with OPT start date is issued AFTER you apply for OPT.}

The sequence is:

\begin{enumerate}
\item You apply for OPT with your school's international office (DSO)
\item DSO issues you a new I-20 with your OPT recommendation and start date
\item You submit Form I-765 to USCIS with that I-20
\item USCIS approves and sends you an EAD card with your OPT dates
\item \textbf{Then} you report employment to SEVP portal
\end{enumerate}

\textbf{Where to find your OPT start date:}

\begin{itemize}
\item \textbf{EAD card} (Employment Authorization Document) - shows your OPT validity period
\item \textbf{OPT I-20} - issued by your DSO when they recommended OPT
\item \textbf{I-797 approval notice} - if you received one from USCIS
\end{itemize}

If you haven't received your EAD card yet, you cannot legally start working and shouldn't be updating employment information in SEVP yet. You must wait until you have your EAD card in hand with your approved OPT dates.

Do you have your EAD card already?
\end{botresponse}

\begin{userprompt}
okay but I think that my offer letter says that 18 Jan 2026 as eligibility start day and my OTP says my start date is 19 Jan 2026. So 18 is Sunday and 19 is Monday and no one works on sunday right ?
\end{userprompt}

\begin{botresponse}
You're correct that January 18, 2026 is a Saturday (not Sunday), and January 19, 2026 is a Sunday.

\textbf{Here's what matters:}

Your OPT start date is \textbf{January 19, 2026} according to your EAD card. This is the earliest date you can legally work on OPT.

\textbf{For the SEVP portal:}

\begin{itemize}
\item You cannot enter January 18, 2026 because it's before your OPT start date
\item Your employment start date must be \textbf{January 19, 2026 or later}
\end{itemize}

\textbf{What to do:}

\begin{itemize}
\item \textbf{Enter January 19, 2026} (your OPT start date) as your employment start date, OR
\item \textbf{Enter the actual first day you physically start working} (likely January 20 or 21 if it's a weekday)
\end{itemize}

The "eligibility start day" on your offer letter likely just means you're eligible to start then, but since your OPT doesn't begin until January 19, that's your legal start date. Most likely you'll actually begin working on Monday, January 20, 2026.

Check with your employer about your actual first work day and use that date (as long as it's on or after January 19).
\end{botresponse}

\begin{userprompt}
no Jan 18 is Sunday
\end{userprompt}

\begin{botresponse}
You're absolutely right—I apologize for the error. January 18, 2026 is Sunday and January 19, 2026 is Monday.

\textbf{For the SEVP portal:}

Since your OPT start date is January 19, 2026 (Monday) and your offer letter says eligibility starts January 18, 2026 (Sunday):

\begin{itemize}
\item You cannot enter January 18 because it's before your OPT start date
\item \textbf{Enter January 19, 2026} as your employment start date
\end{itemize}

This makes sense because:

\begin{itemize}
\item January 19 is your OPT start date (the earliest you can legally work)
\item It's a Monday (a normal work day)
\item The Sunday date on your offer letter is likely just administrative timing
\end{itemize}

Use \textbf{January 19, 2026} on the SEVP portal.
\end{botresponse}

\begin{userprompt}
so does anything needs to be changed in my offer letter ? the date is okay right ?
\end{userprompt}

\begin{botresponse}
No, nothing needs to be changed in your offer letter. The date is fine.

\textbf{Here's why:}

Your offer letter says you're \textbf{eligible} to start January 18, 2026. This is just when you're permitted to begin—it doesn't mean you're required to work that specific Sunday.

Since your OPT legally begins January 19, 2026 (Monday), that's when you'll actually start working. The one-day difference between "eligibility" (Jan 18) and your actual start date (Jan 19) is not a problem.

\textbf{Bottom line:}

\begin{itemize}
\item Offer letter: No changes needed
\item SEVP portal: Enter January 19, 2026 as your employment start date
\item Actual work: You'll start Monday, January 19, 2026
\end{itemize}

The offer letter is perfectly fine as-is.
\end{botresponse}

\begin{userprompt}
do I need to mention my degree when I write what my work duties are in the portal ?
\end{userprompt}

\begin{botresponse}
No, you do not need to mention your degree when writing your work duties in the SEVP portal.

\textbf{What to include:}

\begin{itemize}
\item Just describe what you actually do in the job
\item Focus on tasks and responsibilities
\item Keep it clear and concise
\end{itemize}

\textbf{What NOT to include:}

\begin{itemize}
\item Your degree name
\item Phrases like "using my Applied Data Science degree"
\item Educational background
\end{itemize}

The SEVP system already has your degree information from your I-20. The work duties section is only for describing the job itself.

Use the description I provided earlier—it focuses on your actual responsibilities without mentioning your degree.
\end{botresponse}

\begin{userprompt}
should I mention STAGE center in the duties ?
\end{userprompt}

\begin{botresponse}
No, you don't need to mention STAGE center in the duties description.

\textbf{Why:}

\begin{itemize}
\item You'll already enter the employer name separately (STAGE Center, University of Chicago)
\item The duties section is for describing \textbf{what you do} , not \textbf{where you do it}
\item Keeps the description focused and concise
\end{itemize}

\textbf{Use this version instead:}

I work as part of an interdisciplinary team producing scientific content for theatre and film productions. My responsibilities include testing and comparing AI-based translation tools for content translation and transcription, evaluating speech-to-text software, segmenting transcripts, developing structured instructions to improve AI translation quality, and reviewing outputs for accuracy and consistency. I also assist with video editing and preparation of production materials, applying my data science background to support scientific communication projects.

This is cleaner and avoids redundancy since STAGE center is already listed as your employer.
\end{botresponse}